% FAC-EN-U0062
% Source: fac_body.tex lines 2342--2403.
% Authority: PDF physical pp. 62--63; printed pp. 257--258
% (Chapter III, no. 64 complete).
% Transcription repair FAC-FR-T0036 restores "allons" and four membership signs.

\phantomsection\label{fac:III:3:no:64}
\textbf{64. Comparison of the groups $H^q(M)$ and
$H^q(X,\mathcal{A}(M))$.}

Let $M$ be a graded $S$-module, and let $\mathcal{A}(M)$ be the algebraic
sheaf on $X=\mathbf{P}_r(K)$ constructed from $M$ by the procedure of
\hyperref[fac:III:2:no:57]{no.~57}. We shall compare $C(M)$ with
$C'(\mathfrak{U},\mathcal{A}(M))$, the complex of alternating cochains of
the covering $\mathfrak{U}=\{U_i\}_{i\in I}$ with values in the sheaf
$\mathcal{A}(M)$.

Let $m\in C_k^q(M)$, and let $(i_0,\ldots,i_q)$ be a sequence of $q+1$
elements of $I$. The polynomial $(t_{i_0}\cdots t_{i_q})^k$ plainly
belongs to $S(U_{i_0\cdots i_q})$, with the notation of
\hyperref[fac:III:2:no:57]{no.~57}. It follows
that
\[
\frac{m\langle i_0\cdots i_q\rangle}
     {(t_{i_0}\cdots t_{i_q})^k}
\]
belongs to $M_U$, where $U=U_{i_0\cdots i_q}$, and hence defines a
section of $\mathcal{A}(M)$ over $U_{i_0\cdots i_q}$. As
$(i_0,\ldots,i_q)$ varies, the system formed by these sections is an
alternating $q$-cochain of $\mathfrak{U}$ with values in
$\mathcal{A}(M)$, which we denote by $\iota_k(m)$. It is immediate that
$\iota_k$ commutes with $d$, and that
$\iota_k=\iota_h\circ\rho_k^h$ if $h\geq k$. Passing to the direct limit,
the $\iota_k$ therefore define a homomorphism
\phantomsection\label{fac:III:3:no:64:map:iota}
\[
\iota:C(M)\longrightarrow C'(\mathfrak{U},\mathcal{A}(M))
\]
which commutes with $d$.

\phantomsection\label{fac:III:3:no:64:prop:4}
\textsc{Proposition 4.} \textit{If $M$
\hyperref[fac:III:2:no:56:def:TF]{satisfies condition}}
\hyperref[fac:III:2:no:56:def:TF]{(TF)}\textit{, the homomorphism
$\iota:C(M)\to C'(\mathfrak{U},\mathcal{A}(M))$ is bijective.}

If $M\in\mathcal{C}$, then $M_n=0$ for $n\geq n_0$, whence
$C_k(M)=0$ for $k\geq n_0$ and $C(M)=0$. Since every $S$-module
\hyperref[fac:III:2:no:56:def:TF]{satisfying (TF)} is
$\mathcal{C}$-isomorphic to a module of finite type,
this shows that we may restrict ourselves to the case in which $M$ is of
finite type. We can then find an exact sequence
\[
L^1\longrightarrow L^0\longrightarrow M\longrightarrow0,
\]
where $L^1$ and $L^0$ are free of finite type. By
\hyperref[fac:III:2:no:58:prop:3]{Propositions~3} and
\hyperref[fac:III:2:no:58:prop:5]{5 of no.~58}, the sequence
\phantomsection\label{fac:III:3:no:64:formula:sheaf-exact}
\[
\mathcal{A}(L^1)\longrightarrow\mathcal{A}(L^0)
\longrightarrow\mathcal{A}(M)\longrightarrow0
\]
is an exact sequence of coherent algebraic sheaves. Since the
$U_{i_0\cdots i_q}$ are affine open subsets, the sequence
\phantomsection\label{fac:III:3:no:64:formula:cochain-exact}
\[
C'(\mathfrak{U},\mathcal{A}(L^1))
\longrightarrow C'(\mathfrak{U},\mathcal{A}(L^0))
\longrightarrow C'(\mathfrak{U},\mathcal{A}(M))
\longrightarrow0
\]
is exact (cf.
\hyperref[fac:II:3:no:45:thm:2:cor:2]{no.~45, Corollary~2 to Theorem~2}).
The commutative diagram
\phantomsection\label{fac:III:3:no:64:diagram:iota}
\[
\begin{tikzcd}
C(L^1) \arrow[r] \arrow[d,"\iota"']
  & C(L^0) \arrow[r] \arrow[d,"\iota"']
  & C(M) \arrow[r] \arrow[d,"\iota"'] & 0 \arrow[d] \\
C'(\mathfrak{U},\mathcal{A}(L^1)) \arrow[r]
  & C'(\mathfrak{U},\mathcal{A}(L^0)) \arrow[r]
  & C'(\mathfrak{U},\mathcal{A}(M)) \arrow[r] & 0
\end{tikzcd}
\]
then shows that, if \hyperref[fac:III:3:no:64:prop:4]{the proposition}
holds for $L^1$ and $L^0$, it holds
for $M$. We are therefore reduced to the special case of a free module of
finite type and then, by decomposition into a direct sum, to the case
$M=S(n)$.

In this case $\mathcal{A}(S(n))=\mathcal{O}(n)$. By the very definition
of this sheaf, a section $f_{i_0\cdots i_q}$ of $\mathcal{O}(n)$ over
$U_{i_0\cdots i_q}$ is a regular function on
$V_{i_0}\cap\cdots\cap V_{i_q}$, homogeneous of degree $n$. Since
$V_{i_0}\cap\cdots\cap V_{i_q}$ is the set of points of $K^{r+1}$ at
which $t_{i_0}\cdots t_{i_q}$ is nonzero, there is an integer $k$ such
that
\phantomsection\label{fac:III:3:no:64:formula:fraction}
\[
f_{i_0\cdots i_q}
=\frac{P\langle i_0\cdots i_q\rangle}
       {(t_{i_0}\cdots t_{i_q})^k},
\]
where $P\langle i_0\cdots i_q\rangle$ is a homogeneous polynomial of
degree $n+k(q+1)$, that is, of degree $k(q+1)$ in $S(n)$. Thus every
alternating cochain
$f\in C'(\mathfrak{U},\mathcal{O}(n))$ defines a system
$P\langle i_0\cdots i_q\rangle$ which is an element of $C_k(S(n))$,
and hence a homomorphism
\phantomsection\label{fac:III:3:no:64:map:nu}
\[
\nu:C'(\mathfrak{U},\mathcal{O}(n))\longrightarrow C(S(n)).
\]
Since one immediately verifies that $\iota\circ\nu=1$ and
$\nu\circ\iota=1$, it follows that $\iota$ is bijective, which completes
the proof.

\phantomsection\label{fac:III:3:no:64:cor}
\textsc{Corollary.} \textit{The homomorphism $\iota$ defines an
isomorphism from $H^q(M)$ onto $H^q(X,\mathcal{A}(M))$ for every
$q\geq0$.}

Indeed, we know that
$H'^q(\mathfrak{U},\mathcal{A}(M))=H^q(\mathfrak{U},\mathcal{A}(M))$
(\hyperref[fac:I:3:no:20:prop:2]{no.~20, Proposition~2}), and
$H^q(\mathfrak{U},\mathcal{A}(M))=H^q(X,\mathcal{A}(M))$
(\hyperref[fac:III:1:prop:2]{no.~52, Proposition~2}, which applies since
$\mathcal{A}(M)$ is coherent).

\phantomsection\label{fac:III:3:no:64:remark}
\textsc{Remark.} It is easy to see that
$\iota:C(M)\to C'(\mathfrak{U},\mathcal{A}(M))$ is \textit{injective}
even when $M$ does not
\hyperref[fac:III:2:no:56:def:TF]{satisfy condition (TF)}.
